%% ======================================================================
%% 04a_kinematics.tex -- 3. Model (part 1): assumptions, frames, parameters, kinematics
%% Symbolic verification: scripts_modelo/derive_*.py ; numbers: scripts_modelo/analysis.py
%% ======================================================================
\section{Kinematic and dynamic model}
\label{sec:model}

This section derives the model of the full chain---base $\to$ spherical joint $\to$ octocopter $\to$ plate $\to$ ball---used to size the prototype, design and diagnose its three controllers, and interpret the experiments in Section~\ref{sec:validation}. The derivation proceeds from the bottom up: assumptions, frames, and parameters (Section~\ref{sec:assumptions}); kinematics of the base, the full chain, and the plate actuation (Sections~\ref{sec:basekin}--\ref{sec:linkage}); dynamics of the base, the octocopter, and the ball (Sections~\ref{sec:basedyn}--\ref{sec:ball}); the rationale for using a drone on the ramp and its thrust margin versus battery voltage (Sections~\ref{sec:porquedron}--\ref{sec:bateria}); and finally, the coupled linear model, the controllers, and the resulting fundamental limitations (Sections~\ref{sec:coupling}--\ref{sec:limits_fund}). Every equation for the ball and octocopter dynamics was verified symbolically in SymPy, and all numerical values in this section are computed by the scripts in \fn{scripts_modelo/} in the repository (\ref{app:derivations}).

\subsection{Assumptions, frames, and parameters}
\label{sec:assumptions}

\paragraph{Modeling assumptions}
\begin{enumerate}[label=\textbf{A\arabic*.}, leftmargin=*]
  \item All links are rigid; the compliance of the PLA plates and towers is neglected.
  \item The driven wheels roll without longitudinal or lateral slip; the omnidirectional wheels are passive supports with negligible lateral resistance.
  \item The spherical joint is ideal in roll and pitch; its friction torque $\vect{\tau}_f$ is treated as an unmodeled disturbance. Its yaw is locked by friction and mechanical stops, so the octocopter heading follows the base (confirmed by the yaw cutoffs in Section~\ref{sec:integrated}).
  \item The vertical guide is a prismatic joint with Coulomb and viscous friction $F_g(\dot d)$ and rigid end stops at $d=0$ and $d=d_{\max}$.
  \item Each rotor produces thrust $F_i$ along $+\vect{z}_D$ and a proportional reaction torque; the ESC--motor--propeller assembly responds to its command with a pure time delay followed by a first-order lag $\tau_m$. Thrust at a fixed command scales with the square of the voltage reaching the ESCs (Section~\ref{sec:bateria}).
  \item The ball is a homogeneous sphere rolling without slip; rolling resistance is neglected. The inertia factor $\gamma_b=1/(1+I_b/m_br_b^2)$ is kept general ($\gamma_b=5/7$ for a solid sphere, $3/5$ for a thin spherical shell).
  \item The servos are position servos with a pure time delay and a first-order lag; linkage backlash is treated as a nonlinearity outside the linear model.
  \item For linear analysis, angles, angular rates, and accelerations are small, and the base moves at constant slope except at the ramp transitions.
\end{enumerate}

\paragraph{Frames and notation}
Figure~\ref{fig:frames} shows the coordinate frames. $\{W\}$ is an inertial frame with $\vect z_W$ pointing upward. $\{B\}$ is fixed to the base, with its origin on the ground below the midpoint of the driven (rear) axle, $\vect x_B$ pointing forward (toward the omnidirectional wheels), and $\vect z_B$ normal to the ground. $\{J\}$ has its origin at the center of the spherical joint, at height $h_{J0}+d$ above the ground and distance $\ell_J$ ahead of the driven axle, where $d\in[0,d_{\max}]$ is the displacement along the vertical guide. $\{D\}$ is fixed to the octocopter with its origin at the spherical joint center, $\vect x_D$ toward motor M8, $\vect y_D$ toward motor M7, and $\vect z_D$ upward (the firmware FLU convention). $\{G\}$ is fixed to the octocopter at the center of the plate universal joint, at height $h_P$ along $\vect z_D$. $\{P\}$ is the plate frame, with its origin at the universal joint center; the ball center has coordinates $\vect\rho=[x_b,\,y_b,\,\varepsilon]\T$ in $\{P\}$, where $\varepsilon$ is the height of the ball center above the universal joint center. $\Rot_x,\Rot_y,\Rot_z$ are elementary rotations, $\skw{\vect a}$ is the skew-symmetric matrix of $\vect a$, and orientations use the $z$--$y$--$x$ Euler angles $(\psi,\theta,\phi)$ implemented in the firmware. A vector without a superscript is expressed in the frame specified in the text; $\vect a^{B}$ denotes components in $\{B\}$.

The robot configuration is described by
\begin{equation}
  \vect{q}=\big[\,x,\ y,\ \psi\ \big|\ d\ \big|\ \phi,\ \theta,\ \psi_D\ \big|\ \alpha_x,\ \alpha_y\ \big|\ x_b,\ y_b\,\big]\T\in\mathbb{R}^{11},
  \label{eq:q}
\end{equation}
where $(x,y,\psi)$ is the planar pose of the base in the ground plane, $d$ is the vertical guide displacement, $(\phi,\theta,\psi_D)$ are the roll, pitch, and heading of the octocopter relative to $\{W\}$, $(\alpha_x,\alpha_y)$ are the servo-driven plate tilt angles relative to $\{D\}$, and $(x_b,y_b)$ is the ball position on the plate. The yaw lock in A3 eliminates $\psi_D$ and leaves ten configuration coordinates; the base is subject to one nonholonomic constraint (Section~\ref{sec:basekin}), so the robot has nine velocity degrees of freedom, while the ball's rolling constraints fix its angular velocity in terms of $(\dot x_b,\dot y_b)$ (Section~\ref{sec:ball}). Seven actuation channels drive these nine degrees of freedom (two wheels; the eight rotors, acting as one collective and two torques; two servos), while two coordinates, $(x_b,y_b)$, are unactuated.

\begin{figure}[!htbp]
\centering
\begin{tikzpicture}[scale=1.32, font=\footnotesize, ax/.style={-{Latex[length=1.8mm]}, thick}]
  % ground and world frame
  \fill[pGris!20] (-4.5,-0.32) rectangle (4.5,-0.45);
  \draw[ax] (-4.3,-0.15) -- (-3.6,-0.15) node[above=1pt]{$x_W$};
  \draw[ax] (-4.3,-0.15) -- (-4.3,0.55) node[left]{$z_W$};
  % base: driven (rear) wheel and omni (front) wheel
  \draw[thick, fill=pAzul!10] (-1.6,0.15) rectangle (1.6,0.8);
  \draw[thick, fill=pGris!35] (-1.1,0.16) circle (0.48); \draw[thick, fill=pGris!12] (1.1,0.16) circle (0.48);
  \node at (-1.1,0.16) {driven}; \node at (1.1,0.16) {omni};
  \fill (-1.1,-0.32) circle (0.045);
  \draw[ax, pAzul] (-1.1,-0.32) -- (-0.25,-0.32) node[below right]{$x_B$};
  \node[pAzul] at (-1.88,-0.10) {$\{B\}$};
  % mast offset l_J
  \draw[{Latex}-{Latex}] (-1.1,1.05) -- (-0.4,1.05); \node[above] at (-0.75,1.05) {$\ell_J$};
  \draw[line width=2pt, pGris!85!black] (-0.4,0.8) -- (-0.4,3.0);
  \draw[thick, fill=pGris!22] (-0.6,2.4) rectangle (-0.2,3.0);
  \draw[{Latex}-{Latex}] (-0.85,2.4) -- (-0.85,3.0); \node[left] at (-0.88,2.7) {$d$};
  \draw[{Latex}-{Latex}] (-3.15,-0.32) -- (-3.15,3.1); \node[left] at (-3.20,1.4) {$h_{J0}+d$};
  \draw[dotted] (-3.25,3.1) -- (-0.5,3.1);
  % joint
  \fill (-0.4,3.1) circle (0.09); \node[below left] at (-0.92,3.18) {$\{J\}$};
  \draw[dashed, pGris] (-0.4,3.1) -- (-0.4,7.0);
  % octocopter pitched by theta (all lengths in its own frame)
  \begin{scope}[shift={(-0.4,3.1)}, rotate=-8]
    \draw[line width=1.5pt, pRojo] (-3.1,0.3) -- (3.1,0.3);
    \foreach \xx in {-3.1,-1.7,1.7,3.1} {\draw[thick] (\xx,0.3) -- (\xx,0.5); \draw[thick] (\xx-0.55,0.5) -- (\xx+0.55,0.5);}
    \draw[ax, pRojo] (0,0) -- (1.25,0); \node[pRojo, below] at (1.15,-0.10) {$x_D$ (to M8)};
    \draw[ax, pRojo] (0,0) -- (0,0.75); \node[pRojo] at (0.36,0.62) {$z_D$};
    % centre of mass
    \fill[pRojo!85!black] (0.22,1.05) circle (0.07); \node[right, pRojo!85!black] at (0.30,1.08) {CM};
    \draw[{Latex}-{Latex}] (-0.45,0) -- (-0.45,1.05); \node[left] at (-0.48,0.6) {$h_c$};
    % pillar, universal joint G and plate
    \draw[thick] (0,0.3) -- (0,1.95);
    \fill[pOcre] (0,2.0) circle (0.07); \node[below left, pOcre] at (-0.14,1.94) {$\{G\}$};
    \draw[{Latex}-{Latex}] (2.82,0) -- (2.82,2.0); \node[right] at (2.85,1.0) {$h_P$};
    \draw[dotted] (0.1,2.0) -- (2.88,2.0);
    \begin{scope}[shift={(0,2.0)}, rotate=7]
      \draw[very thick, pOcre, fill=pOcre!18] (-1.6,0.06) rectangle (1.6,0.18);
      \fill[pAzul] (0.75,0.38) circle (0.2);
      \draw[ax, pOcre] (0,-0.04) -- (1.05,-0.04); \node[pOcre, below] at (0.68,-0.06) {$x_P$};
      \draw[pGris, dashed] (0.75,0.38) -- (1.95,0.38); \draw[pGris, dashed] (1.6,0) -- (1.95,0);
      \draw[{Latex}-{Latex}] (1.9,0) -- (1.9,0.38); \node[right] at (1.93,0.19) {$\varepsilon$};
      \draw[{Latex}-{Latex}] (0,0.68) -- (0.75,0.68); \node[above] at (0.37,0.69) {$x_b$};
      \node[pOcre, right, align=left] at (2.30,0.22) {servo tilt $\alpha_y$\\[-1pt]rel.\ to $\{D\}$};
    \end{scope}
    % camera towers (fixed to the drone)
    \draw[thick, pGris] (-2.3,0.3) -- (-0.2,3.6); \draw[thick, pGris] (2.3,0.3) -- (0.2,3.6);
    \draw[thick, fill=pGris!45] (-0.35,3.55) rectangle (0.35,3.85); \node[right] at (0.45,3.75) {camera, fixed to $\{D\}$};
  \end{scope}
  \draw[-{Latex}] (-0.4,6.12) arc (90:82:3.02); \node at (-0.06,6.36) {$\theta$};
\end{tikzpicture}
\caption{Coordinate frames and lengths of the kinematic chain in the $x$--$z$ plane (pitch $\theta$ is shown; roll is analogous about $x$). Frame $\{D\}$ has its origin at the spherical joint center $\{J\}$. The camera towers are fixed to the octocopter, so the camera measures the ball in a frame that tilts with the octocopter but not with the servos.}
\label{fig:frames}
\end{figure}

\paragraph{Parameters}
Table~\ref{tab:params} summarizes the model parameters and their sources. Dimensions come from the CAD assembly of the final prototype (383 parts); masses from weighing; actuator specifications from datasheets; gains from firmware; and three quantities ($\delta_x,\delta_y$, $h_c$, and $k_F$) are identified from flight logs in Section~\ref{sec:ident}.

\begin{small}
\begin{longtable}{@{}>{\raggedright\arraybackslash}p{0.16\textwidth}>{\raggedright\arraybackslash}p{0.43\textwidth}>{\raggedleft\arraybackslash}p{0.35\textwidth}@{}}
\caption{Model parameters. Sources: [CAD] assembly model; [M] measured; [DS] datasheet; [FW] firmware or LabVIEW code; [ID] identified in this work; [E] engineering estimate.}\label{tab:params}\\
\toprule
Symbol & Description & Value [source] \\
\midrule\endfirsthead
\toprule Symbol & Description & Value [source] \\ \midrule\endhead
\multicolumn{3}{@{}l}{\emph{Mobile base}}\\
$r$, $b$, $L$ & wheel radius, track width, wheelbase & \SI{49.3}{\milli\metre}, \SI{0.3794}{\metre}, \SI{0.220}{\metre} [FW, M] \\
$\kappa_e$ & wheel travel per encoder count (1440 counts/rev) & \SI{0.2151}{\milli\metre} [FW, DS] \\
$\tau_s$, $\omega_0$, $i_s$ & TorqueNADO stall torque, no-load speed, and stall current (12\,V, 60:1) & \SI{4.94}{\newton\metre}, \SI{100}{rpm}, \SI{8.7}{\ampere} [DS] \cite{pitsco_torquenado} \\
$M$ & total robot mass with all batteries (remainder \SI{6.8}{\kilo\gram} + module \SI{3.6}{\kilo\gram}) & \SI{10.4}{\kilo\gram} [M] \\
$x_G$, $h_G$ & CM ahead of driven axle and above support plane, from measured tip-over slopes (\SI{21.5}{\degree} and \SI{30.5}{\degree}, module locked) & \SI{0.088}{\metre}, \SI{0.224}{\metre} [M] \\
$\ell_J$, $h_{J0}$, $d_{\max}$ & mast offset, spherical joint height at $d=0$, guide stroke & \SI{0.088}{\metre}, \SI{0.495}{\metre}, \SI{0.140}{\metre} [CAD, M] \\
$C_{rr}$ & rolling resistance coefficient & 0.02 [E] \\
\midrule
\multicolumn{3}{@{}l}{\emph{Aerial module}}\\
$m$, $m_c$ & rotating assembly, carriage & \SI{3.3}{\kilo\gram}, \SI{0.3}{\kilo\gram} [M] \\
$\vect J$ & inertia about spherical joint center & $\diag(0.12;\ 0.12;\ 0.20)$\,\si{\kilo\gram\metre\squared} [CAD] \\
$h_c$ & CM height above spherical joint & $0.14\pm0.06$\,m [ID]; 0.10\,m [CAD] \\
$\delta_x$, $\delta_y$ & horizontal CM eccentricity & $9.5\pm1.8$\,mm, $1.9\pm1.5$\,mm [ID] \\
$L_r$, $h_r$ & rotor radius, rotor plane above spherical joint & \SI{0.350}{\metre}, \SI{0.095}{\metre} [CAD] \\
$(x_i,y_i)$, M1--M4 & rotor positions, clockwise ($s_i=-1$) & $(D,-D)$, $(-D,-D)$, $(-D,D)$, $(D,D)$, $D=L_r/\sqrt2$ [FW] \\
$(x_i,y_i)$, M5--M8 & rotor positions, counterclockwise ($s_i=+1$) & $(0,-L_r)$, $(-L_r,0)$, $(0,L_r)$, $(L_r,0)$ [FW] \\
$k_F$, $u_0$, $n$ & thrust map $F=k_F((u-u_0)/\SI{860}{\micro\second})^n$ & 5.40\,N [FW], 7.17--8.18\,N [ID] at 12.4\,V; \SI{1140}{\micro\second}; 1.6 \\
$c_Q$ & reaction torque per unit thrust & \SI{0.020}{\metre} [FW] \\
$\tau_m$, $T_a$ & motor--ESC lag, actuation delay (50\,Hz ESC frame, IMU filter) & \SI{0.07}{\second} [E], \SI{0.02}{\second} [FW] \\
$C_b$, $R_\Sigma$ & capacity per 3S LiPo; internal + cable + contact resistance per branch & \SI{3.3}{\ampere\hour} [DS]; 41--\SI{68}{\milli\ohm} [E, audit] \\
\midrule
\multicolumn{3}{@{}l}{\emph{Ball and plate}}\\
$h_P$ & universal joint center above spherical joint & \SI{0.104}{\metre} [CAD] \\
$\varepsilon$, $r_b$ & ball center height above universal joint, ball radius & \SI{36.9}{\milli\metre}, \SI{13.5}{\milli\metre} [CAD, M] \\
$r_a$, $\ell_r$, $R_h$ & servo arm, connecting rod, and linkage horn radius & 20, 36, \SI{45.85}{\milli\metre} [CAD] \\
$k_s$ & servo gain (S3003: \SI{45}{\degree} per \SI{400}{\micro\second}) & \SI{0.1125}{\degree\per\micro\second} [DS] \cite{futaba_s3003} \\
$K_\alpha$ & plate tilt per microsecond of pulse width & \SI{0.0491}{\degree\per\micro\second} [CAD, DS] \\
$H_c$ & camera height above plate & \SI{0.20}{\metre} [CAD] \\
\bottomrule
\end{longtable}
\end{small}

\subsection{Differential base kinematics}
\label{sec:basekin}

Let $v$ and $\omega$ be the linear and angular velocities of $\{B\}$, and let $\omega_R,\omega_L$ be the angular velocities of the right and left wheels. By A2, the velocity of the axle midpoint has no lateral component, yielding the Pfaffian (nonholonomic) constraint
\begin{equation}
  \big[-\sin\psi\ \ \cos\psi\ \ 0\big]\,\big[\dot x\ \ \dot y\ \ \dot\psi\big]\T=0,
  \label{eq:pfaff}
\end{equation}
and the kinematic model \cite{siegwart2011,campion1996}
\begin{equation}
  \begin{bmatrix}\dot x\\ \dot y\\ \dot\psi\end{bmatrix}=\begin{bmatrix}\cos\psi&0\\ \sin\psi&0\\ 0&1\end{bmatrix}\begin{bmatrix}v\\ \omega\end{bmatrix},\qquad
  \begin{bmatrix}v\\ \omega\end{bmatrix}=\begin{bmatrix}r/2 & r/2\\ r/b & -r/b\end{bmatrix}\begin{bmatrix}\omega_R\\ \omega_L\end{bmatrix},
  \label{eq:diff}
\end{equation}
whose inverse, used on the myRIO, is $v_{R,L}=r\omega_{R,L}=v\pm\tfrac{b}{2}\omega$, with $b/2=\SI{0.1897}{\metre}$. The instantaneous center of rotation lies on the axle line at distance $R_c=\tfrac{b}{2}\,(\omega_R+\omega_L)/(\omega_R-\omega_L)$ from the midpoint. The two turns specified in the course guide are therefore pivoting about the axle midpoint ($\omega_L=-\omega_R$, $R_c=0$, $\omega=2v_R/b$) and pivoting about one wheel ($\omega_L=0$, $R_c=b/2$, $\omega=v_R/b$). Because the spherical joint lies a distance $\ell_J$ ahead of the axle, it traces a circle of radius $\ell_J=\SI{0.088}{\metre}$ in the first case and $\sqrt{(b/2)^2+\ell_J^2}=\SI{0.209}{\metre}$ in the second; Section~\ref{sec:coupling} shows that the resulting centripetal accelerations are negligible.

\paragraph{Odometry}
The LabVIEW node integrates the encoder increments $\Delta N_{R,L}$ at each sample period using the midpoint rule:
\begin{gather}
  \Delta s_{R,L}=\kappa_e\,\Delta N_{R,L},\qquad \Delta s=\tfrac12(\Delta s_R+\Delta s_L),\qquad \Delta\psi=\frac{\Delta s_R-\Delta s_L}{b},\\
  \begin{bmatrix}x\\y\end{bmatrix}_{k+1}=\begin{bmatrix}x\\y\end{bmatrix}_{k}+\Delta s\begin{bmatrix}\cos(\psi_k+\Delta\psi/2)\\ \sin(\psi_k+\Delta\psi/2)\end{bmatrix},\qquad \psi_{k+1}=\psi_k+\Delta\psi .
  \label{eq:odo}
\end{gather}
For constant wheel velocities within a sample period, the exact (circular-arc) update multiplies the increment by $\operatorname{sinc}(\Delta\psi/2)=1-\Delta\psi^2/24+O(\Delta\psi^4)$, so the truncation error of \eqref{eq:odo} is $\Delta s\,\Delta\psi^2/24$ per step---under \SI{1}{\nano\metre} at the operating speeds ($|\omega|\le\SI{0.5}{\radian\per\second}$, \SI{10}{\milli\second} period). Resolution is one count: $\kappa_e=\SI{0.2151}{\milli\metre}$ in translation and $\kappa_e/b=\SI{0.57}{\milli\radian}$ ($0.032^\circ$) in heading. The dominant odometry errors are therefore systematic \cite{borenstein1996}: a relative difference $\epsilon_d$ between wheel radii produces a heading drift $\epsilon_d S/b$ and a lateral offset $\epsilon_d S^2/(2b)$ over a straight run of length $S$ (\SI{1.5}{\degree} and \SI{13}{\milli\metre} per meter for $\epsilon_d=1\%$), and a relative track-width error $\epsilon_b$ produces a heading error $-\epsilon_b\Theta$ over an in-place turn $\Theta$. Both are calibrated with the 1\,m and 360$^\circ$ tests in Section~\ref{sec:integrated}.

\subsection{Slope embedding and kinematic chain}
\label{sec:chain}

Let the ground be a plane of slope $\beta$ whose direction of steepest ascent is $\vect x_W$. With $(x,y,\psi)$ measured in that plane, the base pose is
\begin{equation}
  \Rot_{WB}=\Rot_y(-\beta)\,\Rot_z(\psi),\qquad \vect p_{B}=\Rot_y(-\beta)\,[x\ \ y\ \ 0]\T ,
  \label{eq:slope}
\end{equation}
and gravity, expressed in $\{B\}$, is
\begin{equation}
  \vect g^{B}=\Rot_{WB}\T\begin{bmatrix}0\\0\\-g\end{bmatrix}=g\begin{bmatrix}-\sin\beta\cos\psi\\ \sin\beta\sin\psi\\ -\cos\beta\end{bmatrix}.
  \label{eq:gB}
\end{equation}
The positions of the spherical joint, universal joint, and ball follow from the chain of homogeneous transformations:
\begin{align}
  \vect p_J&=\vect p_B+\Rot_{WB}\big[\ell_J\ \ 0\ \ h_{J0}+d\big]\T, \label{eq:pJ}\\
  \vect p_G&=\vect p_J+\Rot_{WD}\,h_P\ez,\qquad \Rot_{WD}=\Rot_z(\psi_D)\Rot_y(\theta)\Rot_x(\phi), \label{eq:pG}\\
  \vect p_b&=\vect p_G+\Rot_{WD}\Rot_{DP}\,\vect\rho,\qquad \Rot_{DP}=\Rot_y(\alpha_y)\Rot_x(\alpha_x),\qquad \vect\rho=[x_b\ \ y_b\ \ \varepsilon]\T . \label{eq:pb}
\end{align}
Two properties of \eqref{eq:pG}--\eqref{eq:pb} shape the entire design. First, $(\phi,\theta)$ are \emph{absolute} angles: the octocopter IMU estimates them relative to gravity and its controller drives them to zero. To first order, the absolute tilt of the plate is
\begin{equation}
  \Theta_P\simeq\theta+\alpha_y,\qquad \Phi_P\simeq\phi+\alpha_x ,
  \label{eq:tilt}
\end{equation}
which does not contain the slope $\beta$: the passive spherical joint and self-leveling octocopter decouple ground slope from the plate. If the tray were fixed to the base, the plate would tilt by $\arcsin(\sin\beta\cos\psi)$; on the \SI{10}{\degree} ramp the servos would spend \SI{10}{\degree} of their \SI{\pm10.8}{\degree} stroke (Section~\ref{sec:linkage}) just to level it, and on a \SI{20}{\degree} ramp they could not level it at all. Second, the ball sits at height $h_P+\varepsilon=\SI{0.141}{\metre}$ above the spherical joint, a lever arm that converts octocopter angular accelerations into linear ball accelerations.

\paragraph{Velocity and acceleration propagation}
With $\vect r_{BJ}=[\ell_J,0,h_{J0}+d]\T$, base angular velocity and acceleration $\vect\omega_B,\vect\alpha_B$, and octocopter body-frame angular rates $\vect\omega_D=[p,q,r]\T$, the accelerations of the spherical joint and universal joint are
\begin{align}
  \vect a_J&=\vect a_B+\vect\alpha_B\times\vect r_{BJ}+\vect\omega_B\times(\vect\omega_B\times\vect r_{BJ})+2\,\vect\omega_B\times\dot d\,\ez+\ddot d\,\ez, \label{eq:aJ}\\
  \vect a_G&=\vect a_J+\vect\alpha_D\times h_P\ez+\vect\omega_D\times(\vect\omega_D\times h_P\ez). \label{eq:aG}
\end{align}
For planar motion on a constant slope ($\vect\omega_B=\omega\,\vect z_B$) and $\ddot d=0$, \eqref{eq:aJ} reduces in $\{B\}$ to
\begin{equation}
  \vect a_J^{B}=\big[\dot v-\omega^2\ell_J,\ \ v\omega+\dot\omega\,\ell_J,\ \ 0\big]\T ,
  \label{eq:aJplanar}
\end{equation}
because the axle midpoint acceleration of a nonholonomic base is $[\dot v,\ v\omega,\ 0]\T$. Terms with $\vect\omega_B\times$ and $\vect\alpha_B\times$ in the vertical plane appear only at the ramp transitions, where the base pitches and the spherical joint, \SI{0.5}{\metre} above the ground, experiences the tangential acceleration $(h_{J0}+d)\,\ddot\theta_B$. The octocopter Euler-angle derivatives relate to body angular velocities via
\begin{equation}
  \dot\phi=p+\tan\theta\,(q\sin\phi+r\cos\phi),\quad \dot\theta=q\cos\phi-r\sin\phi,\quad \dot\psi_D=\frac{q\sin\phi+r\cos\phi}{\cos\theta},
  \label{eq:euler_rates}
\end{equation}
which is the integration law of the attitude estimator (\ref{app:derivations}).

\subsection{Plate actuation and ball measurement}
\label{sec:linkage}

\paragraph{Four-bar linkage}
Each plate axis is driven by a four-bar linkage in a vertical plane (Fig.~\ref{fig:ballstab}b): the servo shaft $S$ is the fixed pivot of a crank of length $r_a$, a connecting rod of length $\ell_r$ joins the crank tip to a horn of radius $R_h$ on the universal joint, and the universal joint center $G$ is the second fixed pivot. With $\vect d=G-S=(d_x;\,d_z)=(-25.85;\,36.0)$\,mm, the loop-closure equation $\|S+r_a\vect e(\gamma)-G-R_h\vect e(\alpha)\|=\ell_r$, where $\vect e(\cdot)=[\cos(\cdot),\sin(\cdot)]\T$, becomes the Freudenstein-type equation \cite{freudenstein1955}
\begin{gather}
  A(\gamma)\cos\alpha+B(\gamma)\sin\alpha=C(\gamma), \nonumber\\
  \alpha(\gamma)=\operatorname{atan2}(B,A)+\arccos\frac{C}{\sqrt{A^2+B^2}}, \label{eq:freud}\\
  A=2R_h(r_a\cos\gamma-d_x),\qquad B=2R_h(r_a\sin\gamma-d_z), \nonumber\\
  C=r_a^2+d_x^2+d_z^2+R_h^2-\ell_r^2-2r_a(d_x\cos\gamma+d_z\sin\gamma), \nonumber
\end{gather}
where $\gamma$ and $\alpha$ are the crank and plate angles (on the branch containing the assembled neutral position). Because at neutral the connecting rod is vertical while the crank and horn are horizontal, the small-signal gain is $G_\ell=\dd\alpha/\dd\gamma|_0=r_a/R_h=0.436$. With the servo gain $k_s=\SI{45}{\degree}/\SI{400}{\micro\second}$ \cite{futaba_s3003}, the plate tilts by $K_\alpha=G_\ell k_s=\SI{0.0491}{\degree\per\micro\second}$ ($8.56\times10^{-4}$\,rad/\si{\micro\second}), and the firmware limit of $\pm\SI{220}{\micro\second}$ corresponds to $\pm\SI{10.8}{\degree}$ of plate tilt. Over the full $\pm\SI{60}{\degree}$ servo range, \eqref{eq:freud} spans \SI{-23.0}{\degree} to \SI{21.3}{\degree}, a 7\% asymmetry neglected in the linear model. The backlash measured via servo steps (6/12\,\si{\micro\second} in $x$, 35/15\,\si{\micro\second} in $y$ for positive/negative directions) corresponds to 0.3--\SI{1.7}{\degree} of plate tilt; the firmware adds it to the command with a \SI{30}{\micro\second} ramp.

\paragraph{Camera measurement}
The camera is rigidly attached to the octocopter, so it measures the ball in the plate plane \emph{as seen from} $\{D\}$. A homography $\vect H$ computed from the four plate corners at neutral maps image points to plate millimeters \cite{hartley2004}. Three systematic effects follow from geometry. (i) The ball center lies $r_b$ above the plate surface, so its image projects onto the plate plane with a constant scale factor $H_c/(H_c-r_b)=1.07$, which is absorbed into the loop gain. (ii) A servo tilt $\alpha$ moves the plate relative to the camera and homography; to second order, the reading becomes $x_b\cos\alpha+r_b\sin\alpha$, an error below \SI{2.5}{\milli\metre} at extreme tilt. (iii) Detection runs on the raw image (about 1.04\,px/mm on the plate), and the centroid standard deviation is \SI{0.1}{\milli\metre} for a stationary ball (75\,s log). The Raspberry Pi estimates velocity as the least-squares slope of positions over a window $W=\SI{60}{\milli\second}$ ($N\approx12$ samples at $h=\SI{5.35}{\milli\second}$):
\begin{gather}
  \hat v_k=\sum_{j=0}^{N-1}w_j\,x_{k-j},\qquad w_j=\frac{t_{k-j}-\bar t}{\sum_i (t_{k-i}-\bar t)^2}, \label{eq:lsvel}\\
  \sigma_{\hat v}=\sigma_x\sqrt{\frac{12}{N(N^2-1)}}\;\frac1h\approx\SI{1.6}{\milli\metre\per\second}, \nonumber
\end{gather}
an unbiased FIR differentiator for constant velocity whose low-frequency phase lag matches a delay of $W/2=\SI{30}{\milli\second}$ and whose gain rolls off above $\approx1/W$. Equation~\eqref{eq:lsvel} is used directly, as an FIR filter, in the closed-loop analysis of Section~\ref{sec:ballcontrol}.
